% LaTeX template that I edit and compile on Overleaf

\documentclass[oneside]{article}
\usepackage{essay1}

\usepackage{lmodern}
\usepackage{hyperref}
\usepackage{setspace}

\title{Experience in Formal Education and Industry}
\author{Eric Charles Grasby}
\date{3/16/2026}
\prof{Submitted in partial fulfillment of the application requirements for the Graduate Institute at St. John's College}

\doublespacing

\begin{document}
\maketitle{}
My collegiate education has helped shape the person that I am today. It has allowed me to carve out a career as a software engineer and, additionally, pay for the PhD program that I have chosen to pursue (Computer and Information Science). I want to firmly state that I have never regretted pursuing my existing college education. 

My undergraduate days at UA Little Rock\footnote{\url{https://ualr.edu/}} afforded me the opportunity to diversify my education. I had more chances to take electives than others might have, and I spent many of those elective hours taking courses in the department of rhetoric and writing. This opened a path to the realm of both the art and science of writing that have served me in both my graduate work and in my career as a software developer.

As I transitioned into my first programmer role, I began to really understand what the philosophy of being a developer is rooted in: \textit{problem solving}. If you can think algorithmically (i.e. imagine breaking a problem down into a series of repeatable, manageable steps), then you can apply that thinking to a wide domain of expertise. This is, somewhat, evidenced by the number of programmer roles that I have assumed in various industries for the past decade.

My first career role out of college was in the gilded industry of digital marketing. I spent three years streamlining the process of sending bulk, pre-approved, "junk mail" to hundreds of millions of Americans (talk about an uninspiring job!). Leaving that position behind at the outset of the 2020 pandemic, I managed to land a job in DevOps at a local startup company that performed real-time scam call detection. I viewed this as an attempt at redeeming my proverbial soul after the years of slinging junk mail. 

My time at that startup was tempestuous. At times, I felt like I was cast adrift in that role. Sleepless nights due to production issues led to stressful days. "Free beer" is of little solace to a person that has no opportunity to savor it. Wanting more out of my professional life than the mere "pat on the back" that management (at the time) offered, I chose to, \textit{once again} put myself back on the job market. I later found myself in the service of the IT department for a regional transmission organization (RTO), a company that works to manage the bulk electric system for a corridor of the United States.

\end{document}